% Generated draft: requires mathematical and source review.
\documentclass[11pt]{article}
\usepackage[margin=25mm]{geometry}
\usepackage{fontspec}
\setmainfont{Latin Modern Roman}
\usepackage{amsmath,amssymb,mathtools,unicode-math}
\setmathfont{Latin Modern Math}
\usepackage{enumitem,needspace,xurl,hyperref}
\setlength{\parindent}{0pt}
\setlength{\parskip}{5pt}
\setlength{\emergencystretch}{3em}
\newcommand{\N}{\mathbb{N}}
\newcommand{\Z}{\mathbb{Z}}
\newcommand{\Q}{\mathbb{Q}}
\newcommand{\R}{\mathbb{R}}
\newcommand{\C}{\mathbb{C}}
\newcommand{\F}{\mathbb{F}}
\newcommand{\A}{\mathbb{A}}
\newcommand{\PP}{\mathbb{P}}
\newcommand{\E}{\mathbb{E}}
\newcommand{\eps}{\varepsilon}
\newcommand{\dd}{\,\mathrm{d}}
\begin{document}
\section*{UnsolvedMath ID: 30005657}
\textbf{Source problem code:} catalogue entry
\subsection{Definitions and mathematical statement}
For a noetherian scheme \(X\), let \(D_X=D(\operatorname{Qcoh}(X))\) be the unbounded derived category of quasi-coherent \(\mathcal O_X\)-modules. For a commutative differential graded ring \(R\), let \(D_R=D(R)\) be the unbounded derived category of differential graded \(R\)-modules. The hypotheses on \(R\) are cohomological: \(H^i(R)=0\) for every \(i>0\), and there exists an integer \(a\leq 0\) such that \(H^i(R)=0\) for every \(i<a\). No noetherian hypothesis on \(H^0(R)\) is supplied.

For either category \(D\), a \(t\)-structure is a pair \((\mathcal U,\mathcal V)\) of full subcategories, closed under isomorphisms, such that \(\mathcal U[1]\subseteq\mathcal U\), \(\mathcal V[-1]\subseteq\mathcal V\), \(\operatorname{Hom}_D(\mathcal U,\mathcal V[-1])=0\), and every \(E\in D\) admits a distinguished triangle \(U\longrightarrow E\longrightarrow W\longrightarrow U[1]\) with \(U\in\mathcal U\) and \(W\in\mathcal V[-1]\). An object \(C\in D\) is compact if \(\operatorname{Hom}_D(C,-)\) preserves arbitrary coproducts. We use the convention that a \(t\)-structure is compactly generated if there is a set \(\mathcal S\) of compact objects such that
\[
\mathcal V=\{E\in D:\operatorname{Hom}_D(C,E[j])=0\text{ for every }C\in\mathcal S\text{ and every }j<0\}.
\]
The canonical \(t\)-structure is given by \(D_{\mathrm{can}}^{\leq 0}=\{E:H^i(E)=0\text{ for }i>0\}\) and \(D_{\mathrm{can}}^{\geq 0}=\{E:H^i(E)=0\text{ for }i<0\}\).

\textbf{Classification problem.} For every noetherian scheme \(X\), classify the suitable compactly generated \(t\)-structures on \(D_X\). For every cohomologically nonpositive, cohomologically bounded commutative differential graded ring \(R\), classify the suitable compactly generated \(t\)-structures on \(D_R\). The word ``suitable'' is not defined in the supplied record, and the intended parameter sets and classification correspondences are not fully specified. Thus these are classification requests with an unresolved scope, rather than fully determined bijection statements.

The supplied original excerpt begins a more specific scheme formulation. A full subcategory \(\mathcal A\subseteq D_X\) is called a \(t\)-\(\otimes\)-ideal when
\[
Y\in D_{\mathrm{can}}^{\leq 0},\quad E\in\mathcal A
\quad\Longrightarrow\quad
Y\otimes_{\mathcal O_X}^{\mathbf L}E\in\mathcal A.
\]
This gives a possible tensor compatibility requirement on the aisle \(\mathcal U\). However, the excerpt ends at ``Is there then a bijection between'', before either side of the proposed bijection is supplied. The missing correspondence cannot be reconstructed as an exact source statement from this record.

\textbf{Supplied context and qualifications.} The excerpt says that results in its reference [2] essentially settle the scheme problem for Happel--Reiten--Smal\o{} \(t\)-structures. In particular, it attributes to [2, \S5] a bijection between specialization-closed subsets of \(X\) and torsion pairs \((\mathcal T,\mathcal F)\) in \(\operatorname{Qcoh}(X)\) of finite type, described there as having \(\mathcal T\) generated by a set of coherent sheaves, for which \(\mathcal T\) is a tensor ideal. Here a torsion pair satisfies \(\operatorname{Hom}(\mathcal T,\mathcal F)=0\), and every quasi-coherent sheaf \(M\) fits into an exact sequence \(0\to T\to M\to F\to0\) with \(T\in\mathcal T\), \(F\in\mathcal F\); being a tensor ideal means closure under tensoring with every quasi-coherent sheaf. The excerpt cautions that there can be additional finite-type torsion classes which are not tensor ideals, even for \(X=\mathbb P^1\), citing [2, Examples 5.5 and 6.14]. It refers to [2, Lemma 5.9] for geometric relevance and to Thomason [10] for a classification of compactly generated localizations restricted to localizing tensor ideals. The bibliographic identities of [2] and [10] are absent.

The supplied dated assessment labels the problem partially solved and states that tensor-compatible compactly generated \(t\)-structures on noetherian schemes are classified by Thomason filtrations, while a broader question concerning arbitrary suitable \(t\)-structures and connective differential graded rings remains unresolved. This is supplied curation, not an independently verified theorem or current status determination; its exact hypotheses and its relationship to the truncated formulation require review.
\subsection{Short English statement}
The problem asks for classifications of suitable compactly generated \(t\)-structures on derived categories of quasi-coherent sheaves over noetherian schemes and of modules over cohomologically nonpositive, cohomologically bounded commutative differential graded rings. The meaning of ``suitable'' and the precise proposed bijection are missing. The source discusses tensor compatibility and reports partial progress, but those supplied claims and their hypotheses require human mathematical review.
\subsection{Sources}
\begin{itemize}
\item Supplied local UnsolvedMath registry record, permanent ID \(30005657\), exact source problem code \texttt{catalogue entry}, title ``Compactly Generated t-Structures on Schemes and dg Rings'', in the Oberwolfach Reports Open Problems collection. No individual proposer is supplied; the record gives proposed year 2023.
\item Supplied report citation: \emph{Tensor-Triangular Geometry and Interactions} (2024), Oberwolfach Reports, pp.\ 2303--2358, DOI \texttt{10.4171/owr/2023/40}. \url{https://doi.org/10.4171/owr/2023/40}. The record describes these reports as documenting workshops and open-problem discussions. The supplied citation year 2024 and proposed year 2023 are retained as distinct metadata.
\item Sahoo--Dubey, \emph{Compactly generated tensor t-structures}. Supplied literature link: \url{https://arxiv.org/abs/2204.05015}.
\item Hrbek, \emph{Compactly generated t-structures}. Supplied literature link: \url{https://arxiv.org/abs/1806.00078}.
\item The original excerpt cites [2, \S5], [2, Examples 5.5 and 6.14], [2, Lemma 5.9], and Thomason [10], but supplies no matching bibliography. These labels are preserved without identifying them with the linked literature.
\item Literature assessment date: \texttt{2026-08-22}. The ``partially solved'' status and accompanying classification and unresolved-status claims are supplied dated curation derived from the OWR record, not independently verified here. The supplied record also labels its statement \texttt{corrected\_verified} and reports that the official report was checked directly; this is a source-verification assertion of the supplied record, not a claim of verification for this draft.
\end{itemize}
\end{document}
